\section{Axiomatic Approach}

\begin{definition}[\(\sigma\)-Algebra]
    \label{def:sigma-algebra}%
    Suppose \(\Omega\) is a set, and \(\eventSpace \subseteq \powerset\pqty{\Omega}\) is a collection of subsets of \(\Omega\). We say \(\eventSpace\) is a \emph{\(\sigma\)-algebra}, if the following hold:
    \begin{enumerate}
        \item \(\Omega \in \eventSpace\);
        \item if \(A \in \eventSpace\), then the complement \(A\compl \in \eventSpace\);
        \item if \(\pqty{A_n}_{n \in \NN}\) is a countable collection of sets in \(\eventSpace\) (i.e. \(A_n \in \eventSpace\) for all \(n \in \NN\)), then
        \[
            \bigcup_{n \in \NN} A_n \in \eventSpace.
        \]
    \end{enumerate}
\end{definition}

\begin{definition}[Probability Measure]
    \label{def:probability-measure}%
    Suppose \(\eventSpace\) is a \(\sigma\)-algebra on \(\Omega\). A function
    \[
        \functiondomain{\Prob}{\eventSpace}{\ccintv{0}{1}}
    \]
    is called a \emph{probability measure} if
    \begin{enumerate}
        \item \(\Prob\pqty{\Omega} = 1\);
        \item for any countable disjoint collection \(\pqty{A_n}_{n \in \NN}\) in \(\eventSpace\), we have
        \[
            \Prob\pqty{\bigcup_{n \in \NN} A_n} = \sum_{n \in \NN} \Prob\pqty{A_n}.
        \]
    \end{enumerate}

    The value of \(\Prob\pqty{A}\) for some \(A \in \eventSpace\) is the \emph{probability of \(A\)}.
\end{definition}

\begin{definition}[Probability Space]
    \label{def:probability-space}%
    For a set \(\Omega\), a \(\sigma\)-algebra \(\eventSpace\) on \(\Omega\), and a probability measure \(\Prob\) on \(\eventSpace\) and \(\Omega\), the triple
    \[
        \pqty{\Omega, \eventSpace, \Prob}
    \]
    is a \emph{probability space}.
\end{definition}

\begin{definition}[Outcomes, Events]
    \label{def:outcome-event}%
    The elements of the set \(\Omega\) are called \emph{outcomes}, and the elements of the \(\sigma\)-algebra \(\eventSpace\) are called \emph{events}.
\end{definition}

\begin{remark}
    By convention, if \(\Omega\) is countable, then we take \(\eventSpace\) to be \(\powerset\pqty{\Omega}\), the set of all subsets of \(\Omega\).
\end{remark}

We note that we talk about probability of \emph{events}, rather than \emph{outcomes} -- the probability function is defined from \(\eventSpace\), the \(\sigma\)-algebra.

We may conclude some properties of the probability measure \(\Prob\) from the definition.

\begin{proposition}[Properties of \(\Prob\)]
    \label{prop:properties-prob}%
    \begin{itemize}
        \item \(\Prob\pqty{A\compl} = 1 - \Prob\pqty{A}\).
        \item \(\Prob\pqty{\emptyset} = 0\).
        \item \(\Prob\pqty{A} \leq \Prob\pqty{B}\) if \(A \subseteq B\).
        \item \(\Prob\pqty{A \cup B} = \Prob\pqty{A} + \Prob\pqty{B} - \Prob\pqty{A \cap B}\).
    \end{itemize}
\end{proposition}